\chapter*{WSTĘP}
\addcontentsline{toc}{chapter}{\protect\hspace{0.45cm} WSTĘP}

Współczesny rozwój metod sztucznej inteligencji oraz uczenia maszynowego w~znaczącym stopniu zmienił sposób podejmowania decyzji w~wielu kluczowych obszarach gospodarki, medycyny i~przemysłu. Szczególne miejsce w~tych zastosowaniach zajmują dane tabelaryczne, które stanowią podstawowy format zapisu informacji w~bazach medycznych, systemach bankowych, rejestrach ubezpieczeniowych czy bazach telekomunikacyjnych. Pomimo ciągłego rozwoju zaawansowanych architektur głębokich sieci neuronowych, modelowanie danych tabelarycznych wciąż stawia przed inżynierami spore wyzwania. Wynikają one przede wszystkim z~różnorodności typów cech, zróżnicowania ich rozkładów oraz częstej obecności różnorakich zanieczyszczeń.

W~praktyce inżynierii danych powszechnie wiadomo, że jakość wytrenowanego modelu zależy bezpośrednio od jakości danych, którymi został on zasilony. Dane pozyskiwane z~rzeczywistych środowisk (takich jak bazy kliniczne pacjentów czy systemy rozliczeniowe) bardzo rzadko bywają idealnie kompletne i~czyste. Do najczęściej spotykanych problemów należą: brakujące wartości w~tabelach, obecność skrajnych wartości odstających (będących wynikiem usterek aparatury pomiarowej bądź pomyłek przy wprowadzaniu danych) oraz duże różnice w~skalach poszczególnych zmiennych. Zanieczyszczenia te mogą w~znaczący sposób zaburzać działanie klasyfikatorów, prowadząc do błędnych decyzji i~niskiej trafności predykcji.

Niniejsza praca magisterska poświęcona jest zbadaniu i~usystematyzowaniu wpływu wybranych metod czyszczenia oraz wstępnego przetwarzania danych na skuteczność i~stabilność działania czterech popularnych klasyfikatorów: Naiwnego Klasyfikatora Bayesa, Perceptronu Wielowarstwowego MLP, Lasów Losowych oraz algorytmu XGBoost. Badania przeprowadzono w~warunkach kontrolowanego wprowadzania zanieczyszczeń na pięciu zróżnicowanych zbiorach danych, obejmujących m.in.~układowe zapalenia naczyń, choroby serca, diagnostykę cukrzycy, rezygnacje klientów oraz ocenę wiarygodności kredytowej.

Do przeprowadzenia badań zaimplementowano modularny potok przetwarzania danych w~języku Python, który zabezpiecza cały proces przed wyciekiem informacji (ang.~\textit{data leakage}). Analiza uzyskanych wyników, poparta wykresami oraz testami statystycznymi, pozwoliła na sformułowanie praktycznych wniosków i~rekomendacji dotyczących doboru metod czyszczenia danych w~zależności od wybranego modelu.
